% 06_et.tex - evapotranspiration from the water table, computed by Raven
\section{Add evapotranspiration from the water table}

\step{Give the profiles an extinction depth and run}\label{st:et}
The model's own ET package was removed with \cmd{EVT\_MODEL}. Raven can take ET from the water table itself. An
\cmd{EXTINCTION\_DEPTH} in the profiles adds the package \cmd{EVT\_RAVEN} at the top active cells. Its rate is the
evaporation demand of the HRU that Raven's soil could not meet. The rate falls to zero at the extinction depth below the
cell top. Change \cmd{Creek.rvp}:
\diff{t4_b_rvp.diff}
\term{t4_run_b.txt}
\term{t4_summary_b.txt}
\see \cmd{Exiting Gracefully: Successful Simulation}. The files as they are in this step are in \cmd{\$WS/model/step\_10}.
\why water-table ET takes \V{t4:b:et}~m\textsuperscript{3}/d on average (up to \V{t4:b:etmax} in summer), mostly where the
water table is shallow. The seepage to Raven falls to \V{t4:b:seep}~m\textsuperscript{3}/d.
\begin{trouble}
\emph{``EXTINCTION\_DEPTH > 0 \dots the model keeps its EVT \dots package''}: ET from the model and from Raven would both
act. Take the model's EVT or UZF package over, or set \cmd{EXTINCTION\_DEPTH} to 0.
\end{trouble}
